\section{Overloading Rules}\label{sec:rules}
Now, we formalize the overloading rules informally discussed in Section~\ref{sec:motiv:fgf}
for FGFV.

\subsection{Formalization of Overloading Rules}\label{sec:rules:formal}
\begin{figure}[t]
\small
\begin{tabular}{@{}c@{\hskip0.02\textwidth}c@{}}
\begin{minipage}[t]{.48\textwidth}
\typicallabel{No-Dup-Not-Less}

\newjudge{No Duplicates Rule}{\jnoduplicatesrule{d}{d}}

\infrule[No-Dup-Triv]
  { \name{d}\not=\name{d'} }
  { \jnoduplicatesrule{d}{d'} }
\smallskip

\infrule[No-Dup-Not-Less]
  { \neg\bigtuple{\jle{\dom(d)}{\dom(d')}} }
  { \jnoduplicatesrule{d}{d'} }
\smallskip

\infrule[No-Dup-Not-Gtr]
  { \neg\bigtuple{\jle{\dom(d')}{\dom(d)}} }
  { \jnoduplicatesrule{d}{d'} }

\end{minipage}&%
\begin{minipage}[t]{.5\textwidth}

\newjudge{Meet Rule}{\jmethodmeetrule{d}{d}{\setb{d}}}

\typicallabel{Meet-Third}

\infrule[Meet-Triv]
  { \name{d}\not=\name{d'} }
  { \jmethodmeetrule{d}{d'}{\dontcare} }
\smallskip

\infrule[Meet-Excl]
  { \jle{\bigtuple{\dom(d)\existentialmeet\dom(d')}}{\existstype{\verythinmathspace}{\Bottom}} }
  { \jmethodmeetrule{d}{d'}{\dontcare} }
\smallskip

\infrule[Meet-Less]
  { \jle{\dom(d)}{\dom(d')} }
  { \jmethodmeetrule{d}{d'}{\dontcare} }
\smallskip

\infrule[Meet-Gtr]
  { \jle{\dom(d')}{\dom(d)} }
  { \jmethodmeetrule{d}{d'}{\dontcare} }
\end{minipage}
\end{tabular}

\vspace*{-.5em}
\infrule[Meet-Third]
  { d'''\in\setb{d''}  \quad
    \name{d}=\name{d'}=\name{d'''} \quad
    \jequivalent{\dom(d''')}{\bigtuple{\dom(d)\sqcap\dom(d')}} }
  { \jmethodmeetrule{d}{d'}{\setb{d''}} }

\newjudge{Return Type Rule}{\jreturntyperule{d}{d}}

\typicallabel{Return-Not-Less}

\begin{tabular}{@{}c@{\hskip0.05\textwidth}c@{}}
\begin{minipage}[t]{.45\textwidth}
\infrule[Return-Triv]
  { \name{d}\not=\name{d'} }
  { \jreturntyperule{d}{d'} }
\end{minipage}
&
\begin{minipage}[t]{.5\textwidth}
\infrule[Return-Not-Less]
  { \neg\big(\jle{\dom(d)}{\dom(d')}\big) }
  { \jreturntyperule{d}{d'} }
\end{minipage}
\end{tabular}

\infrule[Return-Test]
  { \arrow(d) = \foralltypeb{\kappa}{\arrowtype{\alpha}{\rho}}
    \qquad
    \rep{\kappa = \lobindup{\rep{\chi}}{P}{\rep{\eta}}}
    \qquad
    \arrow(d') = \foralltypeb{\kappa'}{\arrowtype{\alpha'}{\rho'}}
    \\[3pt]
    \rep{\kappa' = \lobindup{\rep{\chi'}}{Q}{\rep{\eta'}}}
    \qquad
    \mathit{distinct}(\rep{P},\rep{Q})
    \qquad
    \jle{\dom(d)}{\dom(d')}
    \\[3pt]
    \jle{\foralltypeb{\kappa}{\arrowtype{\alpha}{\rho}}}{\foralltype{\bar{\kappa},\bar{\kappa'}}{\arrowtype{\intersectiontype{\alpha}{\alpha'}}{\rho'}}} }
  { \jreturntyperule{d}{d'} }

\vspace*{-1.5em}
\caption{Overloading rules for FGFV}
\label{fig:overloadingrules}
\end{figure}

The overloading rules require that every pair in a set of overloaded method
declarations satisfies all three rules defined in Figure~\ref{fig:overloadingrules}.
To determine whether one method declaration is more specific than the other,
while the informal rules in Section~\ref{sec:motiv:fgf} used the subset relation
of their applicable sets, which is not practically checkable,
we use the subtype relation of their domain types as in \citet{oopsla11}.
Note that even though FGFV additionally supports variance annotations of
trait type parameters, the overloading rules for FGF remain the same
because covariant or contravariant types affect only the well-formedness of
method declarations rather than the validity of overloading.

If two method declarations have different names, they vacuously
satisfy all the rules ({\sc [No-Dup-Triv]}, {\sc [Meet-Triv]}, and
{\sc [Return-Triv]}).  The No Duplicates Rule also checks whether the
domain types of two declarations are different; if the domain type of one
declaration is not a subtype of the domain type of the other, they
satisfy the rule ({\sc [No-Dup-Not-Less]} and {\sc [No-Dup-Not-Gtr]}).
The Meet Rule is satisfied if two domain types exclude
each other ([{\sc Meet-Excl}]), they are in a subtype relation
({\sc [Meet-Less]} and {\sc [Meet-Gtr]}), or the domain type of a third
declaration of the same name is the intersection type of the domain types
of two declarations ({\sc [Meet-Third]}).  The Return Type Rule is
satisfied if two domain types are not in a subtype relation
({\sc [Return-Not-Less]}), or if two arrow types derived from the two declarations are
in a subtype relation ({\sc [Return-Test]}).
% We discuss the domain types
% and arrow types of method declarations and their subtype relations in
% the next subsection.

\subsection{Quantified Types}\label{sec:quant}
\begin{figure}[t]
\small
\newjudge{Well-formed substitution}{ \jobeys{\sigma}{\listofb{P}}{\listofb{\chi}}{\listofb{\eta}} }

\infrule
  { \sigma=\substb{\alpha}{P}
    \andalso
    \rep{\jsubtype[\same{\Delta}]{\same{\sigma}\chi}{\alpha}}
    \andalso
    \rep{\jsubtype[\same{\Delta}]{\alpha}{\same{\sigma}\eta}}
  }
  { \jobeys{\sigma}{\listofb{P}}{\listofb{\chi}}{\listofb{\eta}} }
\vspace*{-1em}
\caption{Well-formed substitution}
\label{fig:wellformedsubst}
\end{figure}

\begin{figure}[t]
\small
\typicallabel{x}
\newjudge{Existential inner subtyping}{\jleinnerusing{\Xi}{\Xi}{\sigma}}
\vspace*{-.5em}

\infrule[E-Sub]
  { \Delta' = \Delta \cup \setb{\lobindup{\rep{\chi}}{P}{\rep{\eta}}} \andalso
    \setb{P} \cap \FV(\rep{\rep{\chi'}},\rep{\rep{\eta'}},{\alpha'}) = \emptyset \andalso
    \sigma = \Bigsubstb{\gamma}{Q} \andalso
	\rep{\gamma\neq\Bottom} \\[3pt]
	\jobeys[\Delta']{\sigma}{\listofb{Q}}{\listofb{\bigtypeunion\setb{\chi'}}}{\listofb{\bigtypeintersection\setb{\eta'}}} \andalso
    \jsubtype[\Delta']{\alpha}{\sigma\alpha'} }
  { \jleinnerusingtight{\existstypeb{\lobindup{\rep{\chi}}{P}{\rep{\eta}}}{\alpha}}{\existstypebigb{\lobindup{\rep{\chi'}}{Q}{\rep{\eta'}}}{\alpha'}}{\sigma} }

\newjudge{Universal inner subtyping}{\jleinnerusing{\Upsilon}{\Upsilon}{\sigma}}
\vspace*{-.5em}

\infrule[U-Sub]
  { \Delta' = \Delta \cup \setb{\lobindup{\rep{\chi'}}{Q}{\rep{\eta'}}} \andalso
    \setb{Q} \cap \FV(\rep{\rep\chi},\rep{\rep\eta},{\alpha}) = \emptyset \andalso
    \sigma = \Bigsubstb{\gamma}{P} \andalso
	\rep{\gamma\neq\Bottom} \\[3pt]
	\jobeys[\Delta']{\sigma}{\listofb{P}}{\listofb{\bigtypeunion\setb{\chi}}}{\listofb{\bigtypeintersection\setb{\eta}}} \andalso
    \jsubtype[\Delta']{\sigma\alpha}{\alpha'} }
  { \jleinnerusingtight{\foralltypebigb{\lobindup{\rep{\chi}}{P}{\rep{\eta}}}{\alpha}}{\foralltypeb{\lobindup{\rep{\chi'}}{Q}{\rep{\eta'}}}{\alpha'}}{\sigma} }

\newjudge{Intersection of existential types}{\Xi\existentialmeet\Xi}
\vspace*{-.5em}
\infrule
  { \mathit{distinct}(\bar{P},\bar{P'}) \andalso
    \setb{P'}\cap\FV(\rep{\rep\chi},\rep{\rep\eta},\alpha)=\setb{P}\cap\FV(\rep{\rep{\chi'}},\rep{\rep{\eta'}},\alpha')=\emptyset }
  { \Bigtuple{\existstypeb{\lobindup{\rep{\chi}}{P}{\rep{\eta}}}{\alpha}} \existentialmeet \Bigtuple{\existstypeb{\lobindup{\rep{\chi'}}{P'}{\rep{\eta'}}}{\alpha'}} = 
    \existstype{\bar{\lobindup{\rep{\chi}}{P}{\rep{\eta}}},\bar{\lobindup{\rep{\chi'}}{P'}{\rep{\eta'}}}}{\intersectiontype{\alpha}{\alpha'}} }

\newjudge{Existential reduction}{\jtreds{\Xi}{\Xi}{\sigma}}
\vspace*{.5em}

\begin{tabular}{@{}c@{\hskip0.05\textwidth}c@{}}
\begin{minipage}[t]{.55\textwidth}
\infrule
  {\jcnequiv{\alpha}{\Bottom}{\mathcal{C}} \andalso \toConstraint{\setb{K}}{\mathcal{C}'} \\
   \junify{\mathcal{C} \wedge \mathcal{C}'}{\sigma}{\mathcal{C}''} \andalso \toBounds{\mathcal{C}''}{\setb{K'}}}
  {\jtreds{\existstypeb{K}{\alpha}}{\existstypeb{K'}{\sigma\alpha}}{\sigma}}
\end{minipage}
&
\begin{minipage}[t]{.4\textwidth}
\infrule
  {\text{otherwise}}
  {\jtreds{\existstypeb{K}{\alpha}}{\existstypeb{K}{\alpha}}{\emptysubst}}
\end{minipage}
\end{tabular}
\bigskip

\newjudge{Universal reduction}{\jtreds{\Upsilon}{\Upsilon}{\sigma}}

\vspace*{-2em}
\hspace*{18em}
\begin{minipage}[t]{.5\textwidth}
\infrule
  { \jtreds{\existstypeb{K}{\alpha}}{\existstypeb{K'}{\alpha'}}{\sigma} }
  { \jtreds{\foralltypeb{K}{\arrowtype{\alpha}{\omega}}}{\foralltypeb{K'}{\arrowtype{\alpha'}{\sigma\omega}}}{\sigma} }
\end{minipage}
\bigskip

\newjudge{Existential subtyping}{\jleusing{\Xi}{\Xi}{\sigma} \andalso \jle{\Xi}{\Xi}}

\begin{tabular}{@{}c@{\hskip0.05\textwidth}c@{}}
\begin{minipage}[t]{.55\textwidth}
\infrule
  { \jtreds{\Xi}{\Xi''}{\dontcare} \andalso \jleinnerusing{\Xi''}{\Xi'}{\sigma} }
  { \jleusing{\Xi}{\Xi'}{\sigma} }
\end{minipage}
&
\begin{minipage}[t]{.4\textwidth}
\infrule
  { \jleusing{\Xi}{\Xi'}{\dontcare} }
  { \jle{\Xi}{\Xi'} }
\end{minipage}
\end{tabular}
\medskip

\newjudge{Universal subtyping}{\jleusing{\Upsilon}{\Upsilon}{\sigma} \andalso \jle{\Upsilon}{\Upsilon}}

\begin{tabular}{@{}c@{\hskip0.05\textwidth}c@{}}
\begin{minipage}[t]{.55\textwidth}
\infrule
  { 
\jtreds{\Upsilon'}{\Upsilon''}{\dontcare}
\andalso 
\jleinnerusing{\Upsilon}{\Upsilon''}{\sigma} 
 }
  { \jleusing{\Upsilon}{\Upsilon'}{\sigma} }
\end{minipage}
&
\begin{minipage}[t]{.4\textwidth}
\infrule
  { \jleusing{\Upsilon}{\Upsilon'}{\dontcare} }
  { \jle{\Upsilon}{\Upsilon'} }
\end{minipage}
\end{tabular}

\smallskip
\newjudge{Constraint solving}{\junify{\C}{\sigma}{\C}}

\smallskip
\newjudge{Conversion of type environment to constraints}{\toConstraint{\Delta}{\C}}

\smallskip
\newjudge{Conversion of constraints to type environment}{\toBounds{\C}{\Delta}}

\vspace*{-.5em}
\caption{Subtype relations of quantified types}
\label{fig:quantified}
\end{figure}

As proposed by \citet{bm97}, we represent domain types and arrow types
of generic method declarations as existential types $\Xi$ and universal
types $\Upsilon$, respectively, collectively called \emph{quantified types}.
For presentation brevity, we omit $C$ from $d = (C, M)$ where a method declaration
$M = \methoddef{m}{\rep{K}}{\rep{x\COLON\alpha}}{\rho}{e}$ is declared in
a class $C$.
The domain type of $d$ is $\dom(d) =
\existstypebigb{K}{\tuple{\bar{\alpha}}}$ and its arrow type is
$\arrow(d) = \foralltypebigb{K}{\arrowtype{\tuple{\bar{\alpha}}}{\rho}}$.
% Because we compare domain types to see whether one generic method
% declaration is more specific than the other, which amounts to see whether
% any instance of the former is more specific than an instance of the latter,
% an existential type $\existstypebigb{K}{\alpha}$ is a subtype of the other
% existential type $\existstypebigb{K'}{\alpha'}$ if there exists
% a substitution $\sigma$ for the type parameters bound in $\rep{K'}$ that makes
% $\sigma\alpha'$ a supertype of $\alpha$.
We compare domain types to decide whether one generic method
declaration is more specific than the other;
an existential type $\existstypebigb{K}{\alpha}$ is a subtype of another
existential type $\existstypebigb{K'}{\alpha'}$ if there exists
a well-formed substitution $\sigma$ for the type parameters bound in $\rep{K'}$ that makes
$\sigma\alpha'$ a supertype of $\alpha$.
A substitution is well-formed if it satisfies the bounds of type parameters
as shown in Figure~\ref{fig:wellformedsubst}.
In contrast, we compare arrow types to see whether the return type
of a more specific generic method declaration can be a subtype of the return
type of a less specific one;
a universal type $\foralltypebigb{K}{\alpha}$ is a subtype of the other
universal type $\foralltypebigb{K'}{\alpha'}$ if there exists
a well-formed substitution $\sigma$ for the type parameters bound in $\rep{K}$ that
makes $\sigma\alpha$ a subtype of $\alpha'$.
Figure~\ref{fig:quantified} formally defines these ``inner'' subtyping
rules for quantified types.  (We assume the usual conventions about
using $\alpha$-conversion to rename bound variables when necessary
to satisfy premises that require variable names to be distinct.)
Finally, the intersection type of two
existential types $\existstypebigb{K}{\alpha}$ and
$\existstypebigb{K'}{\alpha'}$ used in [{\sc Meet-Excl}] and [{\sc Meet-Third}] is
$\exists{\big\llbracket\verythinmathspace \bar{K},\bar{K'} \verythinmathspace\big\rrbracket}{\intersectiontype{\alpha}{\alpha'}}$,
because whenever both $\sigma\alpha$ and $\sigma'\alpha'$ are supertypes
of a type $\tau$ for some $\sigma$ and $\sigma'$,
$(\sigma\compose\sigma')\intersectiontype{\alpha}{\alpha'}$ is a supertype of
$\tau$ as well.

While the subtype relation of domain types is a sound
approximation of the subset relation of applicable sets, it is not
complete; when the applicable set of one declaration is a subset
of that of the other, their domain types may not be in a subtype relation.
Consider the following:
\[
\begin{array}{l}
m\ob{P}(x:{\sf ArrayList}\ob{P})\ \dontcare\andalso
m\ob{Q\extends T}(y:{\sf List}\ob{Q})\ \dontcare\andalso
m\ob{R\extends T}(z:{\sf ArrayList}\ob{R})\ \dontcare
\end{array}
\]
for some trait $T$ where ${\sf ArrayList}$ extends ${\sf List}$ and both
are invariant, and call them $d_1$, $d_2$, and $d_3$ in order.
Even though they are a valid overloading because the applicable set of
$d_3$ is the intersection of those of $d_1$ and $d_2$,
the intersection of the two domain types
$\Xi_{\sqcap}=\existstype{P,Q\extends T}{\intersectiontype{{\sf ArrayList}\ob{P}}{{\sf List}\ob{Q}}}$
is not a subtype of the domain type of $d_3$, 
$\Xi_3=\existstype{R\extends T}{{\sf ArrayList}\ob{R}}$.
However, because ${\sf ArrayList}$ extends ${\sf List}$ and both are
invariant, any subtypes of theirs are exclusive unless they are both
${\sf ArrayList}\ob{S}$ for some type $S$ that is a subtype of $T$.
Thus, $\Xi_{\sqcap}$ denotes the same applicable set as
$\existstype{S\extends T}{{\sf ArrayList}\ob{S}}$, which is
a subtype of $\Xi_3$ using $\subst{R}{S}$.


Therefore, in order to reduce the gap between the subset relation of
applicable sets and the subtype relation of domain types,
we extend the subtype relation of existential types using
\emph{existential reduction}.  As defined in Figure~\ref{fig:quantified},
existential reduction infers a simplified existential type from a given
existential type $\existstype{\rep{K}}{\alpha}$.
As defined in Figure~\ref{fig:constraints}, it first collects
constraints $\C$ that make $\alpha$ different from $\Bottom$:
${\Delta}\turnstile{\alpha}\,{\not\equiv}\,{\Bottom}\Rightarrow{\C}$.
It then tries to find types that instantiate the type parameters bound in
$\rep{K}$ or simply collects their bounds by solving $\C$ and the
constraints on the type parameters; it makes a substitution $\sigma$
for the type parameters that have types to instantiate with and
constraints $\C''$ for the type parameters that have only bounds.
Using the substitution $\sigma$ and type parameter bindings $\rep{K'}$
generated from $\C''$, it produces a simplified existential type
$\existstype{\rep{K'}}{\sigma\alpha}$.  Finally, in the extended subtype
relation, one existential type is a subtype of the other if a reduced
existential type of the former is a subtype of the latter
by the inner subtyping rule as shown in Figure~\ref{fig:quantified}.

\begin{figure}[t]
\small
\begin{tabular}{@{}l@{}l@{}}
\begin{minipage}[t]{.4\textwidth}
\textbf{Constraint grammar} \par
\smallskip
\noindent\(
\begin{array}{@{}lclcl@{}}
\C &::=& \FALSE &\mid& \TRUE \\
  &\mid& T <: T &\mid& T \exc T \\
  &\mid& T \not<: T &\mid& T \nexc T \\
  &\mid& \C \wedge \C &\mid& \C \vee \C \\
  \end{array}
\)
\end{minipage}
&
\hspace*{-1em}
\begin{minipage}[t]{0.6\textwidth}
\newjudge{Constraint negation}{\neg \C = \C}
\noindent\(
\begin{array}{@{}l@{\hspace{1ex}}c@{\hspace{1ex}}ll@{\hspace{1ex}}c@{\hspace{1ex}}l@{}}
\neg(\FALSE) &=& \TRUE & \neg(\TRUE) &=& \FALSE \\
\neg(T \subtypeof T') &=& T \nsubtypeof T' & \neg(T \exc T') &=& T \nexc T'\\
\neg(T \nsubtypeof T') &=& T \subtypeof T' & \neg(T \nexc T') &=& T \exc T' \\
\neg(\C \wedge \C) &=& \neg(\C) \vee \neg(\C) & \neg(\C \vee \C) &=& \neg(\C) \wedge \neg(\C) \\
\end{array}
\)
\end{minipage}\\

\begin{minipage}[t]{0.5\textwidth}
\newjudge{Conditional non-subtyping}{\jnsub{\alpha}{\alpha}{\C}}

\infrule
  { \jsub{\alpha}{\eta}{\C} }
  { \jnsub{\alpha}{\eta}{\neg\C} }

\newjudge{Conditional equivalence}{
\begin{tabular}{l}
\jequiv{\alpha}{\alpha}{\C}\\
\jnequiv{\alpha}{\alpha}{\C}
\end{tabular}
}

\infrule
  { \jsub{\alpha}{\eta}{\C} \andalso \jsub{\eta}{\alpha}{\C'} }
  { \jequiv{\alpha}{\eta}{\C \wedge \C'} }

\medskip
\infrule
  { \jequiv{\alpha}{\eta}{\C} }
  { \jnequiv{\alpha}{\eta}{\neg\C} }

\end{minipage}
&
\begin{minipage}[t]{.5\textwidth}
\newjudge{Conditional variance test}{\jcomp{\alpha}{\alpha}{V}{\C}}

\infrule
  { \jsub{\alpha}{\eta}{\C} }
  { \jcomp{\alpha}{\eta}{\bindco}{\C} }

\infrule
  { \jsub{\eta}{\alpha}{\C} }
  { \jcomp{\alpha}{\eta}{\bindcontra}{\C} }

\infrule
  { \jequiv{\alpha}{\eta}{\C} }
  { \jcomp{\alpha}{\eta}{\bindinv}{\C} }
\newjudge{Converse non-equivalence}{\jcnequiv{\alpha}{\alpha}{\C}}

\infrule
  { \jequiv{\alpha}{\eta}{\C} }
  { \jcnequiv{\alpha}{\eta}{\neg\C} }

\end{minipage}
\end{tabular}

\vspace*{-.5em}
\caption{Constraints. The $\wedge$, $\vee$, and $\neg$ operators form a Boolean algebra.}
\label{fig:constraints}
\end{figure}

Collecting constraints $\C$ that make a type $\alpha$ different from
another type $\alpha'$ requires various relations between types such as
subtype, non-subtype, equivalence, non-equivalence, exclusion, and
non-exclusion relations.  A constraint is inductively defined with
boolean literals, true and false, and simple type relations like subtype,
non-subtype, exclusion, and non-exclusion relations as base cases,
and conjunction and disjunction of constraints as inductive steps.
We also defined the various type relations inductively on the structure
of types.  In addition, we defined all the type relations in a way that
supports complete and disjoint case analysis for types.  The entire type
relations are formally defined in Appendix~\ref{app:condrel}.

%Lastly, the subtype relation of arrow types is used only in \rncite{Return-Test}
Arrow types are compared only in rule \rncite{Return-Test},
to see whether the return type of a more specific
generic method declaration can be a subtype of the return type of a less
specific one.  As the informal rule in Section~\ref{sec:motiv:fgf}
specifies, Return Type Rule checks that for two declarations
$d_1$ and $d_2$ where $d_1$ is more specific than $d_2$, and
for any non-$\Bottom$ type $\gamma$ that $d_1$ is applicable to,
the return type of an instance of $d_2$ that is applicable to
$\gamma$ is a supertype of the return type of some instance of
$d_1$ that is applicable to $\gamma$.
When the arrow types of $d_1$ and $d_2$ are
$\foralltypeb{\kappa}{\arrowtype{\alpha}{\rho}}$ and
$\foralltypeb{\kappa'}{\arrowtype{\alpha'}{\rho'}}$, respectively,
\rncite{Return-Test} checks this by checking whether the arrow type of
$d_1$ is a subtype of
$\foralltype{\bar{\kappa},\bar{\kappa'}}{\arrowtype{\intersectiontype{\alpha}{\alpha'}}{\rho'}}$.
Since $d_1$ is more specific than $d_2$ and is applicable to $\gamma$,
there exist substitutions $\sigma_1$ and $\sigma_2$ that make $d_1$ and
$d_2$ applicable to $\gamma$, respectively, thus both $\sigma_1\alpha$
and $\sigma_2\alpha'$ are supertypes of $\gamma$, therefore so is
$(\sigma_1\compose\sigma_2)\intersectiontype{\alpha}{\alpha'}$.
Thus, if the arrow type of $d_1$ is a subtype of
$\foralltype{\bar{\kappa},\bar{\kappa'}}{\arrowtype{\intersectiontype{\alpha}{\alpha'}}{\rho'}}$
using some substitution $\sigma$, it makes $\sigma\alpha$ a supertype of
$\intersectiontype{\alpha}{\alpha'}$, which makes
$(\sigma_1\compose\sigma_2\compose\sigma)\alpha$ a supertype of $\gamma$,
and $\sigma\rho$ a subtype of $\rho'$.
Hence, $\sigma_1\compose\sigma_2\compose\sigma$ is the
one that makes the instance of $d_1$ that is applicable to $\gamma$
and has the return type that is a subtype of the instance of $d_2$.
In summary, \rncite{Return-Test} ensures that
for any instance of $d_2$ instantiated with $\sigma_2$
that is applicable to $\gamma$, there exists an instance of $d_1$
instantiated with $\sigma_1\compose\sigma_2\compose\sigma$
that is applicable to $\gamma$ and has a return type that is a subtype
of the return type of the instance of $d_2$.

% Similarly for the
In the same manner as for
existential types, we extend the subtype relation of
universal types using \emph{universal reduction} as defined in
Figure~\ref{fig:quantified}.  When [{\sc Return-Test}] checks whether one
universal type is a subtype of the other, the domain type of the latter
is an intersection type.  Thus, in order to utilize the extended
expressivity by existential reduction, universal reduction takes
a universally quantified arrow type $\foralltypeb{K}{\arrowtype{\alpha}{\omega}}$,
simplifies its domain type to $\existstypeb{K'}{\alpha'}$ using
$\sigma$ via existential reduction, and produces a simplified
universal type $\foralltypeb{K'}{\arrowtype{\alpha'}{\sigma\omega}}$.
% Finally, in
In
the extended subtype relation, one universal type is a subtype
of the other if the former is a subtype of a reduced universal type of the
latter by the inner subtyping rule as shown in Figure~\ref{fig:quantified}.


\subsection{Unique Existence of the Most Specific Declaration}\label{sec:uniqthm}
For any method invocation, there always exists a unique most specific
method declaration in its set of applicable method declarations.
\begin{theorem}
\label{thm:mostspecific}
Let $\setb{d}=\visible{\tuple{C,m}}$ for $\jwftype{C}$.\\
\begin{tabular}{l@{~}l}
If &$\applicable{\setb{d}}{\alpha}\neq\emptyset$
for $\jwftype{\alpha}$ and $\alpha\neq\Bottom$,\\
then &
$\left|\mostspecific{\applicable{\setb{d}}{\alpha}}\right|=1$.
\end{tabular}
\end{theorem}
\begin{proof}
See Appendix~\ref{app:proofs}.
\end{proof}
% \begin{proof}[Proof Sketch]
% We first show that the more specific relation of method declarations
% is a partial order.  It is clearly reflexive and transitive.  The No
% Duplicates Rule guarantees anti-symmetricity.  Under the partial order,
% a set of applicable method declarations is a meet semilattice. The Meet
% Rule guarantees the existence of a meet for every nonempty subset of
% the set.  Proving the existence of the most specific method declaration
% is straightforward by structural induction on the meet semilattice.
% \end{proof}

